\documentclass[a4paper,11pt]{article}
\usepackage[utf8]{inputenc}
\usepackage[T1]{fontenc}
\usepackage[french]{babel}
\usepackage{lmodern}
\usepackage{amsmath,amssymb,amsthm}
\usepackage{mathrsfs}
\usepackage{enumitem}
\usepackage{multicol}

\setlist[enumerate]{label=\textbf{\textcolor{Couleur8}{\arabic*.}}, left=1em}

% Si vous voulez aussi modifier les listes enumerate imbriquées :
\setlist[enumerate,2]{label=\textcolor{Couleur8}{\alph*)}}
\setlist[enumerate,3]{label=\textcolor{Couleur8}{\roman*)}}
\setlist[enumerate,4]{label=\textcolor{Couleur8}{\Alph*)}}
\usepackage{graphicx}
\usepackage{xcolor}
\usepackage{tcolorbox}
\usepackage{geometry}
\usepackage{fancyhdr}
\usepackage{titlesec}
\usepackage{cancel}
\usepackage{ifthen}
\usepackage{tikz}
\usetikzlibrary{arrows.meta}
\usepackage{tkz-tab}
\usepackage{eurosym}

% OPTION PROF/ÉLÈVE
% Décommentez UNE des deux lignes suivantes :
\newboolean{prof}\setboolean{prof}{true}   % Version professeur (avec corrections des devoirs)
\newboolean{prof}\setboolean{prof}{true}  % Version élève (sans corrections des devoirs)

\definecolor{Couleur1}{HTML}{4A5D7A} %Th
\definecolor{Couleur2}{HTML}{A5B5D6} %Prop
\definecolor{Couleur3}{HTML}{A8C68A} %Méthode
\definecolor{Couleur6}{HTML}{8A9CC1} %Def
\definecolor{Couleur7}{HTML}{E6B459} %Rq Voc
\definecolor{Couleur8}{HTML}{D36740} %Titres
\definecolor{correction}{HTML}{D36740} % Couleur pour les corrections
\definecolor{coopmathscorr}{HTML}{F15929} % Couleur corrections coopmaths

% Configuration des marges
\geometry{
	top=2cm,
	bottom=2cm,
	inner=1.5cm,
	outer=1.5cm,
}

% Police
\renewcommand{\familydefault}{\sfdefault}

% Compteurs
\newcounter{seancenum}
\newcounter{seancenumD}
\setcounter{seancenum}{1}
\setcounter{seancenumD}{1}

% Configuration des en-têtes et pieds de page
\pagestyle{fancy}
\fancyhf{}
\ifthenelse{\boolean{prof}}{
    \fancyhead[L]{\textcolor{Couleur8}{\textbf{Corrections Automatismes - Classe de seconde - Version Professeur}}}
}{
    \fancyhead[L]{\textcolor{Couleur8}{\textbf{Corrections Devoirs - Séance 17 - Classe de seconde}}}
}
\fancyhead[R]{\textcolor{Couleur8}{\textbf{2025-2026}}}
\fancyfoot[L]{\textcolor{Couleur8}{Mme Bertail et Mme Pinto}}
\fancyfoot[R]{\textcolor{Couleur8}{\thepage}}
\renewcommand{\headrulewidth}{1pt}
\renewcommand{\headrule}{\hbox to\headwidth{\color{Couleur8}\leaders\hrule height \headrulewidth\hfill}}

% Environnements personnalisés
\newtcolorbox{seance}[1]{
    colback=Couleur6!10,
    colframe=Couleur1!65,
    fonttitle=\bfseries\large,
    title=Correction Séance \theseancenum #1,
    left=5mm,
    right=5mm,
    top=3mm,
    bottom=3mm,
    arc=0mm,
    after=\stepcounter{seancenum}
}

\newtcolorbox{seanceD}[1]{
    colback=Couleur6!5,
    colframe=Couleur6!45,
    fonttitle=\bfseries\large,
    title=Correction Devoirs - Séance \theseancenumD #1,
    left=5mm,
    right=5mm,
    top=3mm,
    bottom=3mm,
    arc=0mm,
    after=\stepcounter{seancenumD}
}

\begin{document}

\setcounter{seancenumD}{17}
\newpage
\begin{seanceD}{} %17

\begin{enumerate}
\item $\dfrac{-7x+5}{(-3x-6)(4x-5)} > 0$\\
On commence par chercher les éventuelles valeurs interdites :\\
$(-3x-6)(4x-5) = 0$ si et seulement si $-3x-6 = 0$ ou $4x-5 = 0$.\\
$-3x -6 = 0$ donc $x = -2$\\
$4x -5 = 0$ donc $x = \dfrac{5}{4}$\\
Le quotient est défini sur $\mathbb{R} \smallsetminus \{-2; \dfrac{5}{4}\}$.\\

On peut donc en déduire le tableau de signes suivant :\\

\begin{tikzpicture}[baseline, scale=0.35]
\tkzTabInit[lgt=9,deltacl=0.8,espcl=3.5]{ $x$ / 2.5, $-7x+5$ / 2, $-3x-6$ / 2, $4x-5$ / 2, $\dfrac{-7x+5}{(-3x-6)(4x-5)}$ / 4}{ $-\infty$, $-2$, $\dfrac{5}{7}$, $\dfrac{5}{4}$, $+\infty$}
\tkzTabLine{  , +, t, +, z, -, t, -}
\tkzTabLine{  , +, z, -, t, -, t, -}
\tkzTabLine{  , -, t, -, t, -, z, +}
\tkzTabLine{  , -, d, +, z, -, d, +}
\end{tikzpicture}

L'ensemble de solutions de l'inéquation est $S = \bigg] -2 ; \dfrac{5}{7} \bigg[ \cup \bigg] \dfrac{5}{4}; +\infty \bigg[$.

\item $0{,}25\text{ h} = \dfrac{1}{4} \text{ h} = 15 \text{ min}$.\\
Ainsi, $0{,}25$ heure correspond à ${\color{correction}\boldsymbol{15}}$ {\color{correction}minutes}.

\item L'équation $9x^2-x-5=-5$ s'écrit $9x^2-x=0$.\\
En factorisant le premier membre (facteur commun $x$), on obtient $x(9x-1)=0$.\\
On reconnaît une équation produit nul dont les solutions sont : $0$ et $\dfrac{1}{9}$.\\
$\mathscr{S}={\color{correction}\boldsymbol{\{0;\dfrac{1}{9}\}}}$

\item On applique la propriété des puissances de puissances d'un réel :\\
Soit $n\in \mathbb{N}$, et $p \in \mathbb{N}$, on a : $\left(a^{n}\right)^{p}=a^{np}$\\
$\left(4^{n}\right)^{2}=4^{2n}=\left(4^{2}\right)^{n}={\color{correction}\boldsymbol{16^{n}}}$\\
La bonne réponse est la réponse {\bfseries \color{correction}A}.

\item On note $a$ le montant de l'abonnement mensuel et $p$ le prix d'une séance.\\
D'après le premier ticket : $a + 6 \times p = 44$ soit $a + 6p = 44$.\\
D'après le second ticket : $a + 3 \times p = 26$ soit $a + 3 p = 26$\\
En faisant la différence entre ces deux montants, on obtient :\\
$(a + 6p) - (a + 3p) = 44-26$ soit $3p=18$.\\
Donc le montant pour $1$ séance est : $p = 6$\,\euro{}.\\
On peut alors calculer le montant de l'abonnement mensuel :\\
$a = 44 - 6 \times 6 = 44 - 36 = {\color{correction}\boldsymbol{8\,\,\euro{}}}$.

\end{enumerate}

\end{seanceD}

\end{document}